% ==================== 第7章（重构版） ====================
\chapter{讨论与结论}
\label{ch:conclusion}

\begin{revbox}
\noindent\textbf{本章修订说明。}本章按修订后的证据层级重写讨论与结论：router 是接口，确定分支是控制，AnchorEvolve 是可演化的证据状态，全动作损失感知读出是核心经验贡献。原稿“LLM 锚点不可替代”“融合后验取得最高准确率”等表述不再保留；准确率相对 RF 的 +2.7 点只作为点估计，显著结果是相对旧受限 gate 的损失下降。
\end{revbox}

\section{讨论}

路由实验回答第一个研究问题。DeepSeek Flash 和 GLM 5.3 Flash 分别达到 0.995 和 0.990 的路由准确率，混合效用为 0.403 和 0.419，后者恢复 oracle 的 99.3\%。错路由主要发生在包含未来时态、但答案已经在文件中给出的确定问题；错路由的置信度较高（DeepSeek 上为 .9/.85/.85），说明简单的低置信度人工复核抓不到这些错误。Qwen3-8B 的路由准确率降至 0.867，混合效用 0.263，低于 Always-Anchor 的 0.298。这个对比说明契约识别不是简单的时态匹配：较强基座能够利用查询与截止时点的关系，较小基座则容易把语义边界压缩为表面模式。

\begin{revbox}
\noindent\textbf{Router 的定位。}本文不主张 router 结构创新：它是一个零样本 LLM 接口，用来把契约选择显式化。H1 检验的是“契约选择是否可学”，不是“router 分类器是否优于 Adaptive-RAG/Self-RAG 等已有路由”。混合效用的提升来自分支契约的改进（527 条已验证子集上 .389→.437），而不是 router 本身的重新设计。若论文把 router 当第一贡献，审稿人只会看到增量；把它写成接口，才能把创新让给契约形式化、演化状态和读出层。
\end{revbox}

DocFinQA 和 FinDVer 回答第二个研究问题，但它们是\textbf{控制分支}的结果。DeepSeek 上，义务检索分支取得 0.457 精确匹配、平均 10.69 段，固定 16 段阅读器取得 0.390、平均 16 段；GLM 上分别为 0.497 对 0.480，平均 10.31 段对 16 段。FinDVer 上三个基座分别为 0.701、0.753 和 0.649，均高于固定 16 段的 0.660、0.743 和 0.636。四个具名 RAG 适配在 DocFinQA 上均低于义务检索分支，分支的答案覆盖率 0.770 和 0.810 也高于固定预算。

\begin{revbox}
\noindent\textbf{控制分支的边界。}确定分支不是本文的新检索算法；主实验使用确定性的词项义务覆盖，类型化 metadata 过滤是设计目标而非已实现性质。成本口径也要诚实：DeepSeek 上最终分支相对固定 16 段只省 2.3\% 的总 token，且 planner B=16 在准确率和 token 上都支配最终分支；GLM 上最终分支省 29\% token。FinDVer 上 GLM/Qwen 的宏 F1 低于固定 16 段。因此 H2 只支持“控制分支足以让 router 的另一半可信”，不支持“新检索算法一致占优”。
\end{revbox}

收入指引实验回答第三个研究问题，也是本次修订的核心。旧系统（受限 gate）的 Fusion 决策损失为 3.604；全动作损失感知读出在同一后验和同一损失下把它降到 2.863（DeepSeek）、2.956（GLM）和 2.837（Qwen），三个配对区间都排除零。RF 上也从 3.454/3.460/3.454 降到 2.881/2.885/2.881。该读出是严格改进：原 gate 只在点答案与最高密度 $1-\alpha$ 可信集之间选择，新 gate 枚举全部点标签、7 个非空裁决集和弃权。

\begin{revbox}
\noindent\textbf{公平基线。}把所有具名基线放到同一个全动作 gate 下，Fusion 的点估计低于全部基线；Qwen 上显著，DeepSeek/GLM 上相对最强基线不显著（DeepSeek 对 Self-RAG 差 $-0.185$，CI $[-0.581,0.211]$；GLM 对 FinDVer 差 $-0.265$，CI $[-0.699,0.177]$）。原稿“显著低于所有基线”只成立于“baseline 只能点答案、FinPlanRAG 可以用集合/弃权”的不对等设定；修订版不再使用这个比较。
\end{revbox}

文本锚点与时间演化给出\textbf{增量但不是主因}的结果。新 prompt 把 LLM 单独准确率从 .445 提到 .502；新锚点 + 演化 + 演化先验 RF 的最终系统达到 Acc .551、Bal .489、Macro .503、Loss 2.890，相对旧系统损失差 $-0.714$，CI $[-1.057,-0.352]$，显著；相对演化先验 RF 诊断（.555/.503/.508/2.828）的损失差为 $+0.062$，CI $[-0.141,0.273]$，不显著。准确率相对旧系统 +2.7 点，但 McNemar $p\approx0.69$，只能作为点估计。

\begin{revbox}
\noindent\textbf{演化消融。}新锚点 static 2.908→evolved 2.855；旧锚点 static 2.899→evolved 2.793；方向一致但幅度小，227 条上不显著。演化先验 RF 把测试准确率从 .520 提到 .555（+3.5 点，McNemar $p\approx0.40$），Bal/Macro 分别 .476→.503、.473→.508；这些也是点估计。结论应写成“时间演化和演化先验提供了一致的点估计改善，显著增益来自全动作读出”，而不是“演化是主因”。
\end{revbox}

\section{局限与未来研究}

当前证据边界主要来自五点。第一，新文本锚点只跑了 DeepSeek；GLM/Qwen 只验证了全动作读出（旧锚点），需要对应 API key/本地模型才能完成三基座的文本锚点实验。第二，收入指引标签来自 XBRL 自动核验，不是人工金标准；227 条测试中 33 条 below，类别不平衡使平衡准确率/宏 F1 的点估计不稳定；46 条纯文本标签的敏感性结果仍需报告。第三，公司名已脱敏，但 LLM 预训练记忆仍可能污染 2022--2025 的结果；需要用 cutoff 早于 2022 的模型或合成公司做对照。第四，“可审计”是设计属性，尚未做人工审计实验；锚点来源、资格理由和信念变化目前只是可检查的接口，不是已测量的用户收益。第五，全动作裁决集覆盖率约 .86，低于 $1-\alpha$ 集的 .99；这是给定损失下的最优权衡，但不同机构换损失参数后动作分布会改变，需要敏感性分析。

\begin{revbox}
\noindent\textbf{控制分支与实现边界。}确定分支的主结果仍使用 regex 词项覆盖和固定 chunk 时间戳，不是第 4 章的类型化 planner；DocFinQA 只跑 300/7,437 样本；planner B=16 在 DeepSeek 上支配最终分支；FinDVer 的 GLM/Qwen 宏 F1 回退。这些是控制分支的边界，不应被包装成主贡献。类型化 planner、全测试集、dense/hybrid 检索和预算对齐是确定分支的后续工作，不应挤占 AnchorEvolve/readout 的主线。
\end{revbox}

后续研究将围绕五个方向展开。第一，做 LLM 污染检验：公司名脱敏之外的 cutoff 早于 2022 模型、合成公司、或“是否知道结果”的对照。第二，把时间演化扩展到多轮/滚动 origin，并在更多前瞻性事件（生产里程碑、产能投放、监管审批、合同履约）上评估，而不是只做收入指引三桶。第三，补人工审计实验：有/无锚点 provenance 时，审计错误发现率/时间/信心是否改善；如果“可审计”要继续作为贡献，这是必须的。第四，补简单 router 基线（关键词/时点规则、TF-IDF/逻辑回归、统一灵活输出接口的 LLM），把 H1 从“LLM 能分”收紧为“契约路由比简单规则/统一接口更好”。第五，做损失参数、$\alpha$、$w$、$\rho$、$p_e$、硬锚点阈值的敏感性网格，并报告 coverage-size 曲线与 conformal/阈值 gate 对照。

\section{研究结论}

本文把金融混合查询的回答形式选择表述为查询契约的查询级决策，并以义务检索和 AnchorEvolve 分别实现闭合答案契约和前瞻性判断契约。Router 是这一契约选择器的简单接口，不是本文的主要算法创新。确定契约由义务检索控制分支实现；前瞻契约由 AnchorEvolve 实现：从截止时点 8-K 原文抽取带来源的 L2 锚点，经五类资格判定、时间演化的类型可靠性/先验和信念更新得到可审计证据状态，再由全动作损失感知读出转换为点答案、裁决集或弃权。

\begin{revbox}
\noindent 三项贡献的证据层级是：\textbf{契约级形式化}是概念贡献，router 只是接口；\textbf{可演化的锚点证据状态}是状态贡献，其价值在审计、时点正确和作为决策输入，而不是单独点预测；\textbf{全动作损失感知读出}是核心经验贡献，在 DeepSeek、GLM、Qwen 三个基座上显著降低决策损失（3.604→2.863、3.648→2.956、3.577→2.837），并把基线比较放到同一个 gate 下。DeepSeek 演化版最终系统达到 .551/.489/.503/2.890，相对旧系统显著降损（$-0.714$，CI $[-1.057,-0.352]$）；演化先验 RF 诊断达到 .555/.503/.508/2.828。准确率相对 RF 的 +2.7 点不显著，论文不把它写成显著优势。
\end{revbox}

从金融实践看，框架把答案、证据和风险放在同一条审计链上。分析师对确定问题获得带证据编号的闭合答案，对前瞻性断言获得点判断、损失最优裁决集或弃权，并可以查看锚点来源、资格理由和信念变化。机构可以按自身风险偏好替换损失参数和校准后端，而不改变契约选择与证据状态的结构。生产里程碑、产能投放、监管审批和合同履约等任务与收入指引具有相似的时点结构，本文结果为这些场景提供可复用的接口。

本文的结论有明确边界。收入指引实验采用三桶结果空间；文本锚点目前只在 DeepSeek 上完成；标签来自 XBRL 自动核验；227 条测试使准确率/平衡准确率的点估计不稳定；全动作裁决集覆盖率低于 $1-\alpha$ 集；可审计性尚未做人工实验。后续工作将围绕污染检验、多轮演化、人工审计、简单 router 基线和损失参数敏感性展开。由于契约选择、锚点状态和读出彼此分离，这些改进可以独立评估，而不需要改变 FinPlanRAG 的整体接口。
